% ---
% RESUMOS
% ---

% resumo em português
\setlength{\absparsep}{18pt} % ajusta o espaçamento dos parágrafos do resumo
\begin{resumo}
O crescimento exponencial do \textit{streaming} de vídeo impõe desafios críticos às Redes de Distribuição de Conteúdo (CDNs), especialmente quanto ao balanceamento eficiente de carga entre servidores heterogêneos.
%
Estratégias convencionais como \textit{Round-Robin} e Seleção Aleatória ignoram o estado real dos recursos dos servidores, resultando em sobrecarga de nós limitados e degradação da Qualidade de Experiência (QoE).
%
Este trabalho propõe, implementa e valida um algoritmo probabilístico de balanceamento de carga baseado em monitoramento de memória RAM e taxa de leitura de disco, incorporando mecanismos de interpretação de dados desatualizados adaptados do método \textit{Interpreting Stale Load Information}.
%
A solução foi implementada em arquitetura completa com servidores MPEG-DASH, balanceador em \textit{Rust}, comunicação via MQTT e ambiente virtualizado com \textit{Docker}.
%
Experimentos em quatro cenários — variando concorrência (100/1000 usuários) e heterogeneidade de disponibilidade de recursos por servidor — demonstraram reduções de latência de até 85,7\% e eliminação de travamentos em ambientes heterogêneos sob alta carga, validando a superioridade do algoritmo proposto.

 \textbf{Palavras-chave}: 
 %
 Balanceamento de Carga.
 %
 Sistemas Distribuídos.
 %
 Distribuição de Conteúdo.
 %
 Monitoramento de Memória.
 %
 \textit{Streaming} de Vídeo.
 %
 CDN (\textit{Content Delivery Network}).
 %
 \textit{Cache} de Servidores.
 %
 Arquitetura de Sistemas.
 %
 QoE (\textit{Quality of Experience}).
 %  
 Redes de Computadores.
 %
 Heurísticas de Balanceamento.
\end{resumo}
